\section{Introduction}\label{sec:introduction}
Graph Neural Networks (GNNs) are a popular choice for learning tasks on graphs.
Examples include molecules~\cite{gilmer2017mpnn}, road networks~\cite{derrow2021eta},
or social and citation networks~\cite{kipf2017gcn}.
The power of GNNs stems from their ability to encode complex, discrete objects
into dense, semantic vectors serving as the representation of a graph.
They do so by iteratively propagating node features across neighbors, potentially taking
edge features into account.
In a final global pooling step, the node information is fused into a graph-level embedding
representing the graph.
However, recovering the graph input from the embedding poses several challenges, as the
vector only semantically encodes information like graph size, type distributions, and connectivity.

Traditionally, encoding-decoding tasks are formulated under the umbrella term of \red{autoencoders}, where a
neural network learns to map a high-dimensional object to a lower-dimensional code while being
able to map back the code to the input representation.
By mapping an input to a low-dimensional latent representation and training a decoder to reconstruct it,
they extract the most salient features of this input.
There exists a large variety of autoencoders, including the standard autoencoders~\cite{hinton2006autoencoder},
denoising autoencoders~\cite{vincent2008dae}, sparse autoencoders~\cite{ng2011sparseae},
masked autoencoders~\cite{he2022mae}, or variational autoencoders~\cite{vaes}.
They have been applied for representation learning of different modalities like
images~\cite{hinton2006autoencoder,vincent2008dae}, text~\cite{bowman2016generating}, and sound~\cite{zeghidour2021soundstream},
which have well-defined input-output mappings in terms of their representation.

In general, mapping back a graph embedding to a graph representation is challenging for the following reasons.
Besides their discrete and variable-sized nature, which also hold for text, it is their combinatorial
and unordered nature that prevents direct, end-to-end optimization of the input-output mapping.
The fundamental difficulty is related to graph isomorphism, i.e., by permuting the node order, we do not
change the object we represent.
As such, a graph with $n$ nodes can be represented by $n!$ unique but equivalent representations, which
makes gradient descent on the neural network output ill-posed as the loss is not unique for a particular prediction.
As such, autoencoders for graphs are a relatively understudied subfield as the optimization of graphs
as neural network outputs poses unique challenges.
Therefore, the first variational graph autoencoders~\cite{vgae} only decoded the adjacency between
node embeddings rather than the full graph embedding.
However, this does not unlock the full potential as only a graph-level representation is truly meaningful, for example,
to interpolate between variable-sized graphs or generate new ones by modifying the latent code.
As a consequence, a few graph-level methods have been proposed so far, all of which use some form of
node matching~\cite{winter2021pigvae,any2graph2024,grale2025} internally.


Unfortunately, while direct optimization of the output is straightforward and computationally efficient,
it must inherently deal with the problem of graph isomorphism.
In this work, we show that we can avoid dealing with graph isomorphism by framing the problem
as a sequential search procedure over graphs, which removes the dependence on node orderings completely.
By deciding at each step whether a particular graph fragment lies on a path towards the target, we are
able to build latent-to-graph decoders that significantly outperform the state-of-the-art in terms of
generating the correct graph from the latent.

\red{
    Explain why "auto-encoding" graphs may not make a lot of sense since the majority aims to
    solve an ill-posed problem by learning to generate a canonical order.
}
